\chapter{Background: Python's Semantics}
\label{ch:python}

This chapter reviews the Python language constructs that \Deppy{} has
to interpret.  It is shorter than the dependent-type-theory chapter
because most readers are at least familiar with Python; we focus on
the corners that the verifier has to model carefully.

We use ``Python'' to mean CPython 3.10+ unless we specifically need
to discuss earlier versions or alternative interpreters.  The
sidecar pipeline is implemented against Python 3.10 minimum (the
\texttt{match}/\texttt{case} statement is a frequent dependency).

\section{Duck typing and structural subtyping}
\label{sec:duck}

The Python community has long used \emph{duck typing}: ``if it walks
like a duck and quacks like a duck \dots''  Operationally, a
function does not insist that its argument inhabit a particular
nominal type --- it just calls methods on the argument and
\emph{expects them to work}.

\begin{lstlisting}[language=psdl, caption={Duck typing in idiomatic
Python.}]
def total_lengths(items):
    return sum(len(x) for x in items)

# Works on lists, tuples, generators, dicts, strings, custom objects:
total_lengths([[1, 2], (3, 4, 5), "abc"])    # → 8
\end{lstlisting}

\texttt{total\_lengths} works for any iterable whose elements
respond to \texttt{len()}.  Standard \emph{nominal} subtyping
(``\texttt{List} extends \texttt{Iterable}'') would require every
such carrier to declare itself a subclass of some abstract iterable.
Duck typing skips this: the argument just has to \emph{look right}
at call time.

PEP 544 \citep{PEP544} formalises this with \emph{Protocols}:
structurally-typed records of method signatures.  A class implements
\verb|Sized| if it has \verb|__len__| with the right signature, full
stop --- no \verb|class C(Sized)| declaration needed.

\subsection{Why this matters for verification}

The implication for verification is significant.  A property such as
\begin{quote}
``\texttt{total\_lengths} returns a non-negative \texttt{int}.''
\end{quote}
is true \emph{for every input that the function does not reject}.
The set of acceptable inputs is open --- any future class that
implements \texttt{\_\_len\_\_} returning \texttt{int $\ge 0$}
qualifies.  A nominal-types-only verifier cannot say anything about
this property without enumerating all candidate types.

The cubical-duck-type view makes this explicit: a class $C$ that
implements a protocol $P$ \emph{inhabits the same fibre} of the
protocol-fibration as every other class that implements $P$.  The
proof of a protocol-level property descends through the fibration.
We discuss the formal model in \Cref{ch:cubical}; the kernel term is
\verb|DuckPath|.

\section{Method resolution order (MRO)}
\label{sec:mro}

When you call \verb|x.method()|, Python looks up \verb|method| on
\verb|type(x)|, then on \verb|type(x)|'s bases, and so on, in a
linearised MRO order computed by C3 linearisation.

\begin{lstlisting}[language=psdl]
class A:           pass
class B(A):        pass
class C(A):        pass
class D(B, C):     pass
D.__mro__
# (D, B, C, A, object)
\end{lstlisting}

The MRO is \emph{the} static piece of information \Deppy{} uses to
resolve dispatched methods.  When the verifier sees
\texttt{x.distance(y)} on \texttt{x : Point}, it consults
\verb|Point.__mro__| to find the resolved \verb|distance|, reads its
source via \verb|inspect.getsource|, and translates that source.

Multiple inheritance complicates this.  When \verb|distance| is
defined in two ancestors, the C3 linearisation determines which
wins.  \Deppy{} respects the linearisation but does not, in the
current implementation, prove anything about the fact that C3 was
used (it is taken on the user's authority that CPython implements
C3 correctly).

\section{Exceptions and the $\Diamond$-modality}
\label{sec:exceptions}

Python exceptions are first-class control-flow.  Every function
either \emph{returns a value} (the \emph{value fibre}) or \emph{raises
an exception} (the \emph{exception fibre}).  These two fibres are
disjoint --- a function call's outcome is exactly one of them.

The cubical-duck-type model formalises this as a $\Diamond$-modality
on types:
\begin{itemize}[leftmargin=1.7em]
  \item $\Box A$ --- the type of guaranteed-pure expressions of type
    $A$.  No exception, no I/O, no mutation.
  \item $\Diamond A$ --- the type of \emph{maybe-effectful}
    expressions whose value, when one is produced, is of type $A$.
\end{itemize}

A \texttt{try / except} block is a $\Diamond$-eliminator that breaks
into the two fibres:

\begin{lstlisting}[language=psdl, caption={\texttt{try / except} as
$\Diamond$-elimination.}]
try:                                   # ◇ branch
    result = risky_operation(x)
except ValueError:                     # exception fibre
    result = fallback(x)
\end{lstlisting}

In \Deppy{}'s kernel: \texttt{EffectWitness(effect="value", proof=$P_v$)}
plus \texttt{ConditionalEffectWitness(precondition="raises ValueError",
effect="exception", proof=$P_e$)}.

\subsection{Why the modality matters}

A property like ``\texttt{Point.distance(p, q) >= 0}'' is meaningful
\emph{only on the value fibre}.  On the exception fibre, the function
does not return a value, so the property is vacuously not violated.
The \PSDL{} \texttt{else}-branch typically discharges the exception
fibre with a \texttt{raise NotImplementedError} that the kernel
records as a \texttt{Structural} vacuous claim.

The soundness gate (\Cref{ch:soundness}) flags \texttt{Structural}
leaves; in the current implementation, vacuity claims on the
exception fibre are \emph{permitted} but counted, and
\texttt{certify\_or\_die} accepts them only when no
counterexample-search outcome contradicts them.

\section{Monkey-patching}
\label{sec:monkey}

Python permits runtime mutation of classes:
\begin{lstlisting}[language=psdl]
import sympy.geometry as G
G.Point.distance = my_faster_distance     # monkey-patch
\end{lstlisting}

After this assignment every \texttt{Point} instance, including ones
created before the patch, dispatches \verb|distance| to
\verb|my_faster_distance|.  Any property previously verified about
\verb|G.Point.distance| now claims something about
\verb|my_faster_distance|.

\Deppy{} models monkey-patching as a \emph{path} between class
versions:
\begin{equation*}
  \text{C} =_{\mathsf{PyClass}} \text{C[m := v]}
\end{equation*}
which holds \emph{only if} the \emph{behavioural contract} of
\texttt{C.m} survives the patch.  Kernel term: \texttt{Patch}.

The honest position is that we do not verify monkey-patches at the
moment.  The kernel \emph{has} the term, but the pipeline does not
detect monkey-patching at run time --- the user is implicitly
responsible for not patching a class while a verified certificate is
in scope.  A future iteration of the pipeline could compare the
\verb|inspect.getsource| of \verb|C.m| against a recorded SHA-256
(the \verb|expects_sha256:| field on \verb|verify| blocks already
exists, see \Cref{ch:deppy}); the gate currently only checks the
hash, not the patch.

\section{Descriptors, properties, and \texttt{\_\_getattr\_\_}}
\label{sec:descriptors}

Python has \emph{three} attribute-lookup protocols, in priority order:
\begin{enumerate}[leftmargin=1.7em]
  \item \emph{Data descriptors} on \verb|type(x)|
    (e.g.\ \verb|@property|).
  \item Instance dictionary entries.
  \item Non-data descriptors and class attributes.
  \item \verb|__getattr__| on \verb|type(x)| as a fallback.
\end{enumerate}

The \verb|@property| decorator turns a method into a data
descriptor: \verb|circle.radius| calls the getter even though it
looks like attribute access.  The \verb|verify segment_length_code|
block in \texttt{sympy\_geometry.deppy} verifies a property of
\verb|Segment.length| which is a \verb|@property|; the verifier has
to call \verb|method.fget| explicitly to get the underlying
function.  See \verb|sidecar_verify._process_one|, lines around 190.

The \verb|__getattr__| fallback is more subtle.  It is only called
when normal lookup fails, so a class that defines
\verb|__getattr__| effectively has an \emph{open} method set.  In
\Deppy{}, \verb|__getattr__|-based proxies are modelled as
fibered lookups (kernel term \verb|FiberedLookup|): the resolved
attribute is the second projection of a $\Sigma$-type
$\Sigma\,(t : \mathit{Target}).\, \mathsf{methods}(t)$, and the
kernel records the predicate under which the projection inhabits the
protocol fibre.

\section{Metaclasses}
\label{sec:metaclasses}

A metaclass is a class whose instances are themselves classes.
\verb|type| is the default metaclass; many libraries (including
\Sympy) define custom metaclasses for class-registration tricks.

\Deppy{}'s kernel records metaclass-defined hierarchies as
fibrations:
\begin{lstlisting}[language=psdl]
Fiber.from_metaclass(MetaC)
\end{lstlisting}
returns a \verb|Fiber| proof term whose \verb|type_branches|
enumerate registered subclasses
(\verb|deppy/core/kernel.py|).  This is again \emph{advisory}; the
verify-block pipeline does not currently consult metaclass-derived
fibrations except when the user explicitly requests fibration
descent in the \verb|fibration:| field of a \verb|verify| block.

\section{The Python AST}
\label{sec:ast}

The verifier ingests Python source as an \verb|ast.Module| --- the
same AST the CPython compiler walks.  Every \verb|verify| block's
target is resolved to a \verb|FunctionDef| whose body is then
recursively translated to either:
\begin{itemize}[leftmargin=1.7em]
  \item a \Lean{} expression, via
    \verb|deppy.lean.body_translation|;
  \item a kernel \verb|ProofTerm|, via
    \verb|deppy.proofs.psdl| applied to a \verb|psdl_proof| block.
\end{itemize}

The AST nodes that \PSDL{} maps directly to proof-term
constructors are listed in \Cref{ch:psdl}; the body-translation
mapping (relevant when a verify block does not have an explicit
\verb|psdl_proof| but does want its body translated to \Lean{}) is
in \Cref{ch:lean}.

We use \verb|ast.parse(src, mode="exec")| for module-level
parsing and \verb|ast.parse(src, mode="eval")| for property
expressions.  The \verb|ast.unparse| inverse is used heavily in the
emitter to produce \Lean{}-side text.

\section{Notes on dynamic features}
\label{sec:dynamic}

The verifier does not, in general, model:

\begin{itemize}[leftmargin=1.7em]
  \item \verb|exec| / \verb|eval|.
  \item \verb|globals()| / \verb|setattr| modulo runtime conditionals.
  \item Dynamically-generated classes via \verb|type(name, bases, dict)|.
  \item C-extension boundaries (\Sympy{}'s \verb|gmpy2| backend, etc.).
  \item Threads, asyncio scheduling order, GIL semantics.
\end{itemize}

When the body translator hits one of these (most commonly a
C-extension fall-through where \verb|inspect.getsource| fails), it
returns \verb|sorry| with a note.  The verify block records this in
its outcome (\verb|VerifyOutcome.translation_sorry_count > 0|) and
the certificate flags it.

\section{Summary}
\label{sec:ch3-summary}

The Python features the verifier explicitly models:
\begin{itemize}[leftmargin=1.7em]
  \item Duck typing, structurally typed via PEP 544 protocols
    \citep{PEP544}, encoded as protocol fibrations and
    \verb|DuckPath|s.
  \item C3-linearised MRO, taken on faith from CPython, used to
    resolve dispatched method bodies.
  \item Exceptions as the exception fibre of a $\Diamond$-modality;
    \texttt{try/except} as $\Diamond$-elimination.
  \item Monkey-patching as a \verb|Patch| path; not currently
    auto-detected at runtime.
  \item Descriptors, especially \verb|@property|, handled by
    \verb|method.fget| dereferencing.
  \item \verb|__getattr__| proxies via \verb|FiberedLookup|.
  \item Metaclass-derived fibrations via \verb|Fiber.from_metaclass|.
  \item AST-level body translation via \verb|deppy.lean.body_translation|.
\end{itemize}

The dynamic features outside this set fall into the trust surface
described in \Cref{ch:soundness}.  The next chapter formalises the
cubical-duck-type metatheory that ties together the type theory of
\Cref{ch:dependent} with the Python features of this chapter.
